\documentclass[12pt,a4paper]{report}

\usepackage[utf8]{inputenc}
\usepackage[T1]{fontenc}
\usepackage[french]{babel}
\usepackage{geometry}
\usepackage{array}
\usepackage{longtable}
\usepackage{hyperref}

\geometry{margin=2.5cm}

\title{Rapport sur le pipeline de l'application DocuAI}
\author{Projet PFA - Extraction intelligente de documents}
\date{}

\begin{document}

\maketitle

\section*{Introduction}

Ce rapport présente uniquement le pipeline de fonctionnement de l'application DocuAI. Il décrit les étapes suivies depuis l'importation d'un document jusqu'à l'obtention d'un résultat JSON structuré, exploitable dans l'interface web.

L'application est conçue pour traiter plusieurs types de documents : factures STEG, analyses médicales, factures fournisseurs et tickets de caisse. Le pipeline combine le prétraitement d'image, l'OCR local, Docling, PaddleOCR, Qwen local via Ollama et, si nécessaire, l'API Gemini.

\section*{Pipeline global de l'application}

Le pipeline global correspond à l'enchaînement complet des traitements réalisés par l'application.

\begin{center}
\begin{tabular}{c}
Document PDF ou image \\
$\downarrow$ \\
Importation depuis l'interface web \\
$\downarrow$ \\
Validation du fichier \\
$\downarrow$ \\
Prétraitement du document \\
$\downarrow$ \\
Détection du type de document \\
$\downarrow$ \\
Extraction OCR ou IA \\
$\downarrow$ \\
Normalisation des champs en JSON \\
$\downarrow$ \\
Validation métier et score qualité \\
$\downarrow$ \\
Stockage dans l'historique \\
$\downarrow$ \\
Affichage dans l'application \\
\end{tabular}
\end{center}

\section*{Étapes détaillées du pipeline}

\subsection*{1. Importation du document}

L'utilisateur importe un document depuis la page d'extraction de l'application. Les formats acceptés sont principalement :

\begin{itemize}
    \item PDF ;
    \item JPG ;
    \item PNG ;
    \item TIFF ;
    \item WEBP ;
    \item BMP.
\end{itemize}

L'utilisateur peut choisir le type du document ou laisser l'application en mode automatique.

\subsection*{2. Validation du fichier}

Après l'importation, le backend vérifie que le fichier est exploitable. Cette étape contrôle :

\begin{itemize}
    \item le format du fichier ;
    \item la taille maximale autorisée ;
    \item la présence réelle du fichier ;
    \item la lisibilité du document ;
    \item la validité des pages dans le cas d'un PDF.
\end{itemize}

Si le fichier est invalide, l'application retourne une erreur claire à l'utilisateur.

\subsection*{3. Prétraitement du document}

Le prétraitement permet d'améliorer la qualité visuelle du document avant l'extraction. Il est réalisé principalement avec OpenCV.

Les opérations appliquées sont :

\begin{itemize}
    \item correction de l'orientation EXIF ;
    \item redimensionnement de l'image ;
    \item correction de perspective ;
    \item correction de l'inclinaison ;
    \item amélioration du contraste ;
    \item amélioration de la netteté ;
    \item préparation du document pour l'OCR.
\end{itemize}

Cette étape est importante car la qualité de l'image influence directement la qualité de l'OCR et donc la fiabilité des champs extraits.

\subsection*{4. Détection du type de document}

Lorsque le mode automatique est sélectionné, l'application essaie d'identifier le type du document. Cette détection repose sur :

\begin{itemize}
    \item le nom du fichier ;
    \item les mots-clés détectés dans le document ;
    \item le texte extrait rapidement par OCR ;
    \item la structure générale du document.
\end{itemize}

Les types reconnus sont :

\begin{itemize}
    \item facture STEG ;
    \item analyse médicale ;
    \item facture fournisseur ;
    \item ticket de caisse.
\end{itemize}

\subsection*{5. Choix de la méthode d'extraction}

L'application propose plusieurs méthodes d'extraction.

\begin{longtable}{|p{4cm}|p{10cm}|}
\hline
\textbf{Méthode} & \textbf{Rôle} \\
\hline
Pipeline local & Méthode principale basée sur Docling, PaddleOCR et Qwen local via Ollama. \\
\hline
OCR local classique & Méthode de secours utilisant l'OCR sans raisonnement avancé par modèle de langage. \\
\hline
API Gemini & Méthode externe utilisée lorsque l'utilisateur sélectionne Gemini ou pour certains documents complexes. \\
\hline
\end{longtable}

\subsection*{6. Pipeline local Docling + PaddleOCR + Qwen}

Le pipeline local est le pipeline principal de l'application. Il permet de traiter les documents localement, sans envoyer les données vers un service externe.

Son fonctionnement est le suivant :

\begin{center}
\begin{tabular}{c}
Document prétraité \\
$\downarrow$ \\
Docling convertit le document en texte ou Markdown \\
$\downarrow$ \\
Contrôle de qualité du texte extrait \\
$\downarrow$ \\
Si le texte est insuffisant : fallback PaddleOCR \\
$\downarrow$ \\
Si PaddleOCR échoue : fallback Tesseract OCR \\
$\downarrow$ \\
Envoi du texte vers Qwen local via Ollama \\
$\downarrow$ \\
Extraction des champs selon un schéma JSON \\
\end{tabular}
\end{center}

Dans ce pipeline :

\begin{itemize}
    \item Docling sert à convertir le document en contenu structuré ;
    \item PaddleOCR sert à reconnaître le texte lorsque Docling n'est pas suffisant ;
    \item Tesseract OCR sert de solution de secours rapide ;
    \item Qwen local transforme le texte extrait en JSON structuré ;
    \item Ollama exécute le modèle Qwen localement.
\end{itemize}

\subsection*{7. Pipeline Gemini}

Lorsque la méthode Gemini est sélectionnée, le document est traité avec l'API Gemini. Cette méthode peut être utilisée pour comparer les résultats avec le pipeline local ou pour traiter certains documents complexes.

Le résultat retourné par Gemini est ensuite normalisé afin de respecter la même structure JSON que les autres méthodes.

\subsection*{8. Extraction des champs}

Les champs extraits dépendent du type de document.

\begin{longtable}{|p{4cm}|p{10cm}|}
\hline
\textbf{Type de document} & \textbf{Champs principaux extraits} \\
\hline
Facture STEG & Référence client, numéro compteur, date facture, montant à payer, date limite de paiement. \\
\hline
Analyse médicale & Laboratoire, patient, code patient, date de naissance, sexe, numéro dossier, numéro d'examen, dates, médecin demandeur, résultats des examens. \\
\hline
Ticket de caisse & Magasin, date, heure, numéro ticket, articles, total, mode de paiement. \\
\hline
Facture fournisseur & Fournisseur, client, numéro facture, date facture, articles, taxes, total TTC. \\
\hline
\end{longtable}

\subsection*{9. Normalisation JSON}

Après l'extraction, les données sont converties dans un format JSON standardisé. Ce format permet à l'application d'afficher, stocker et exporter les résultats de manière homogène.

Exemple simplifié :

\begin{verbatim}
{
  "document_type": "steg_invoice",
  "reference": "123456789",
  "numero_compteur": "456789",
  "date_facture": "2025-05-12",
  "montant_a_payer": "185.750"
}
\end{verbatim}

\subsection*{10. Validation métier}

L'application vérifie ensuite si les champs importants sont présents. Par exemple :

\begin{itemize}
    \item une facture STEG doit contenir la référence, le numéro compteur, la date facture et le montant à payer ;
    \item une analyse médicale doit contenir le patient et les résultats des examens ;
    \item un ticket de caisse doit contenir le magasin et le total ;
    \item une facture fournisseur doit contenir le numéro facture et le total.
\end{itemize}

Si des champs importants sont absents, l'application ajoute des avertissements dans le résultat JSON.

\subsection*{11. Score qualité}

Un score qualité est calculé afin d'évaluer la fiabilité du résultat extrait. Ce score dépend principalement du nombre de champs importants correctement détectés.

Dans l'interface, ce score permet d'indiquer si le résultat est :

\begin{itemize}
    \item fiable ;
    \item correct ;
    \item à vérifier.
\end{itemize}

\subsection*{12. Stockage et affichage}

À la fin du traitement, l'application enregistre :

\begin{itemize}
    \item le document source ;
    \item le résultat JSON ;
    \item le type du document ;
    \item la méthode utilisée ;
    \item le statut du traitement ;
    \item les avertissements ;
    \item le score qualité.
\end{itemize}

Ces informations sont ensuite affichées dans les pages Documents, Historiques et Tableau de bord. L'utilisateur peut consulter les champs extraits, ouvrir le détail du document et exporter un rapport PDF.

\section*{Conclusion}

Le pipeline de l'application DocuAI permet de transformer un document brut en données structurées et exploitables. Il combine le prétraitement d'image, l'OCR local, Docling, PaddleOCR, Qwen local via Ollama et l'API Gemini.

Cette architecture hybride permet d'améliorer la fiabilité de l'extraction, notamment lorsque les documents sont bruités, inclinés, flous ou de structures différentes.

\end{document}
